% sections/knowledge_graph.tex — Reference Knowledge Graph Construction
\section{Reference Knowledge Graph Construction}
\label{sec:knowledge_graph}

This section describes how \system{} constructs reference knowledge
graphs from the generated enterprise data, encoding all three knowledge
layers (\Cref{def:three_layers}) into graph structures suitable for
GNN-based audit analytics.

% ─────────────────────────────────────────────────────────────────────
\subsection{Graph-Theoretic Foundation}
\label{sec:kg_foundation}

Double-entry bookkeeping has a natural directed-graph
representation~\citep{ijiri1965generalized}: each distinct GL account
forms a node, credit entries produce outgoing edges, and debit entries
produce incoming edges.  Following~\citet{ivertowski2024hardware}, we
adopt a temporal formulation:
\begin{equation}
G(t_0, t_1) = \bigl(A(t_0, t_1),\; E(t_0, t_1)\bigr)
\label{eq:temporal_graph}
\end{equation}
where $A(t_0,t_1)$ is the set of account nodes with monetary state over
$[t_0, t_1]$, and each edge $E_i(t_0,t_1) = \sum_{n=1}^{|\Omega_{A_i}|}
\omega_n$ aggregates all transactions $\omega_n$ affecting node $A_i$
in that period.  This temporal definition enables system-theoretic
analyses such as cross-correlation and auto-correlation on account flows.

Complementing accounting networks constructed from real data,
\system{}'s graphs are generated \emph{forward}: every edge has known provenance, eliminating the
combinatorial matching ambiguity of \Cref{thm:combinatorial}.

% ─────────────────────────────────────────────────────────────────────
\subsection{Entity Type Coverage}
\label{sec:kg_entities}

The reference knowledge graph encompasses 78+ entity types across
the full enterprise data stack, organized by domain:

\begin{table}[t]
\centering
\caption{Entity type coverage in the reference knowledge graph,
  organized by enterprise domain.}
\label{tab:entity_types}
\small
\begin{tabular}{lcp{8cm}}
\toprule
\textbf{Domain} & \textbf{Types} & \textbf{Representative Entities} \\
\midrule
Accounting & 6 & JournalEntry, ChartOfAccounts, ACDOCA,
  AccountBalance, TrialBalance, FiscalPeriod \\
Master Data & 6 & Vendor, Customer, Material, FixedAsset,
  Employee, EntityRegistry \\
Document Flow & 10 & PurchaseOrder, GoodsReceipt, VendorInvoice,
  Payment, SalesOrder, Delivery, CustomerInvoice,
  CustomerReceipt \\
Sourcing (S2C) & 9 & SourcingProject, RfxEvent, SupplierBid,
  ProcurementContract, CatalogItem, SupplierScorecard \\
Financial Reporting & 6 & FinancialStatement, ManagementKpi,
  Budget, ConsolidationSchedule, OperatingSegment \\
HR/Payroll & 6 & PayrollRun, TimeEntry, ExpenseReport,
  BenefitEnrollment, DefinedBenefitPlan, StockGrant \\
Manufacturing & 6 & ProductionOrder, RoutingOperation,
  QualityInspection, CycleCount, BomComponent \\
Tax & 8 & TaxJurisdiction, TaxCode, TaxReturn,
  TaxProvision, DeferredTaxRollforward, UncertainTaxPosition \\
Treasury & 6 & CashPosition, CashForecast, CashPool,
  HedgingInstrument, DebtInstrument, DebtCovenant \\
ESG & 8 & EmissionRecord, EnergyConsumption, WaterUsage,
  WasteRecord, WorkforceDiversityMetric, SafetyIncident \\
Project & 5 & Project, ProjectCostLine, EarnedValueMetric,
  ChangeOrder, ProjectMilestone \\
Audit & 8 & EngagementLetter, AuditOpinion, KeyAuditMatter,
  ComponentAuditor, GroupAuditPlan, SamplingPlan \\
Controls & 4 & InternalControl, ControlMapping, SoD,
  CosoComponent \\
\bottomrule
\end{tabular}
\end{table}

% ─────────────────────────────────────────────────────────────────────
\subsection{Relationship Edge Types}
\label{sec:kg_edges}

Entities are connected by 39+ typed relationship edges with explicit
cardinality constraints:

\begin{itemize}
  \item \textbf{Transactional edges:} Debit/credit flows between
    accounts (one-to-many), document chain references
    (PO$\to$GR$\to$Invoice$\to$Payment), payment allocations.
  \item \textbf{Hierarchical edges:} Chart of accounts parent/child,
    cost center hierarchy, organizational reporting lines, corporate
    group structure.
  \item \textbf{Master data edges:} Vendor-to-material linkages,
    customer-to-segment assignments, employee-to-department.
  \item \textbf{Control edges:} Approval relationships with threshold
    attributes, segregation-of-duties constraint edges, COSO
    principle-to-control mappings.
  \item \textbf{Standards edges:} GAAP rule applicability, ISA
    requirement mappings, SOX assessment linkages.
  \item \textbf{Cross-process edges:} Inventory P2P-O2C links
    (GoodsReceipt $\to$ Delivery), payment-bank reconciliation,
    intercompany bilateral matching.
\end{itemize}

Each edge carries a type discriminator, cardinality constraint
(one-to-one, one-to-many, many-to-many), and domain-specific attributes
(amounts, dates, status flags).

% ─────────────────────────────────────────────────────────────────────
\subsection{Three-Layer Graph Encoding}
\label{sec:kg_encoding}

The reference knowledge graph explicitly encodes all three knowledge
layers as graph properties:

\paragraph{Layer 1 (Structural).}
Node existence, edge connectivity, document chain integrity, master
data relationships.  The graph builder constructs three complementary
subgraphs:

\begin{itemize}
  \item \textbf{Transaction graph:} Nodes are accounts/entities; edges
    are transactions with amount, date, and process-type features.
  \item \textbf{Approval network:} Nodes are users; edges are approval
    relationships with threshold attributes.  SoD conflicts are
    detectable as structural motifs.
  \item \textbf{Entity relationship graph:} Nodes are legal entities,
    vendors, customers; edges represent ownership, control, and
    supply-chain relationships.
\end{itemize}

\paragraph{Layer 2 (Statistical).}
Node and edge properties carry distributional metadata: the
\texttt{ToNodeProperties} trait maps domain models to seven
property-value types (\texttt{String}, \texttt{Decimal},
\texttt{Integer}, \texttt{Boolean}, \texttt{Date}, \texttt{DateTime},
\texttt{List}), ensuring lossless round-tripping through graph
serialization.  Amount distributions, temporal patterns, and
correlation structures are recoverable from the graph properties.

\paragraph{Layer 3 (Normative).}
Control nodes encode COSO component/principle mappings, maturity
levels, and scope assignments.  Standards nodes encode GAAP/IFRS rule
applicability and audit procedure requirements.  Anomaly labels are
attached as edge/node attributes with full provenance (type, category,
difficulty, severity, scheme membership).

% ─────────────────────────────────────────────────────────────────────
\subsection{Hypergraph Extensions}
\label{sec:kg_hypergraph}

Some enterprise relationships are inherently multi-entity.  A hypergraph
builder extends the model to capture these:

\begin{itemize}
  \item \textbf{Three-way match hyperedge:} Links a purchase order,
    goods receipt, and vendor invoice simultaneously.
  \item \textbf{Intercompany elimination:} Links the matched pair of IC
    transactions with their consolidation entry.
  \item \textbf{Period-close:} Links accrual entries, depreciation
    postings, and the resulting financial statement line items.
\end{itemize}

% ─────────────────────────────────────────────────────────────────────
\subsection{Export Formats}
\label{sec:kg_export}

The reference knowledge graph is exported in three formats optimized for
different downstream use cases:

\begin{itemize}
  \item \textbf{PyTorch Geometric (\texttt{.pt}):} Heterogeneous graph
    tensors with typed nodes and edges, directly loadable for GNN
    training.  Node features are scaled and normalized; edge indices
    use COO format.
  \item \textbf{Neo4j (CSV + Cypher):} Node and relationship CSV files
    with Cypher import scripts for graph database analysis, Cypher
    queries, and visualization.
  \item \textbf{DGL (heterograph):} Typed heterogeneous graph with
    node/edge feature dictionaries for Deep Graph Library research.
\end{itemize}

% ─────────────────────────────────────────────────────────────────────
\subsection{Process Mining (OCEL 2.0)}
\label{sec:kg_ocel}

The OCPM module generates Object-Centric Event Logs conforming to the
OCEL~2.0 standard~\citep{ghahfarokhi2021ocel}.  Events are linked to
multiple typed objects via many-to-many relationships.  Lifecycle state
machines for \texttt{PurchaseOrder}, \texttt{SalesOrder}, and
\texttt{VendorInvoice} objects define probabilistic transitions,
producing realistic variant distributions.

Multi-object correlation events---three-way match, payment allocation,
and bank reconciliation---link objects across process boundaries.
Resource pool workload modeling assigns events to users via configurable
strategies (round-robin, least-busy, skill-based).  Enriched events
carry \texttt{from\_state}, \texttt{to\_state},
\texttt{resource\_workload}, and \texttt{correlation\_id} attributes
for conformance checking and resource-perspective analysis.

The OCEL~2.0 logs are fully traceable to the reference knowledge
graph---every event corresponds to a known structural, statistical,
and normative context, enabling process mining research with ground
truth.

% ─────────────────────────────────────────────────────────────────────
\subsection{Ground Truth Properties of the Reference Graph}
\label{sec:kg_ground_truth}

By construction, the reference knowledge graph satisfies the
completeness property of \Cref{thm:completeness}.  Concretely, for any
query over the graph:

\begin{enumerate}
  \item \textbf{Node classification:} Every entity's type, attributes,
    and anomaly status are known.  This provides ground-truth labels for
    node classification tasks in GNN research.
  \item \textbf{Edge prediction:} Every relationship is generated with
    known provenance.  Missing edges in a subgraph can be identified
    against the full graph, enabling link prediction benchmarks.
  \item \textbf{Anomaly detection:} Every injected anomaly is labeled
    with type, difficulty, severity, and causal chain.  Detection
    algorithms can be evaluated against this ground truth across all
    five knowledge dimensions (\Cref{def:dimensions}).
  \item \textbf{Community structure:} Vendor clusters, customer
    segments, intercompany groups, and departmental boundaries are known,
    enabling evaluation of community detection algorithms.
  \item \textbf{Temporal evolution:} The graph can be sliced by period,
    enabling temporal GNN research on time-evolving audit knowledge.
\end{enumerate}
